\section{Two Nuisance Variables Larger Than the Effect}\label{sec:geometry}
Neither of the variables in this section is interesting in itself. Both are reported because
each is larger than several of the backbone differences that frozen-encoder comparisons are run
to establish, and neither is routinely reported at all.

\subsection{How the crop is fitted to the input}
A pedestrian crop is roughly $64\times128$; a ViT input is square. Something has to reconcile
them, and the two conventions in common use disagree about what. \texttt{timm}'s default
evaluation transform resizes the short side and centre-crops, which on such an image retains
roughly the middle 45\% of the body and discards the head and the feet. The
\texttt{torchhub} path resizes to the target directly, distorting the aspect ratio and keeping
everything.

Measured on identical CLIP ViT-B/16 weights at identical input size, with geometry as the only
variable, the choice moves probed Market-1501 mAP from \clipcropMarket{} to \clipMarket{}. It
also moves the head's training accuracy over the source identities from $0.7073$ to $0.9319$,
which is the more diagnostic of the two: under the centre-crop convention an angular margin
cannot separate the training identities in CLIP's frozen space, and under the other convention
it nearly can. A conclusion of the form ``an angular margin fails in CLIP's representation''
--- which an earlier version of this study drew --- was substantially a statement about a
resize transform.

The consequence for the ablation of \cref{sec:ablation} is direct. The teachers reach the bench
through \texttt{timm} and the student through \texttt{torchhub}. Run on each backend's default,
the teachers would have carried a handicap the student never paid, worth up to a large fraction
of their probed score, and \cref{tab:ablation} would have reported the opposite conclusion with
the same confidence. This is why the resize geometry is a hashed field of the encoder
specification in our substrate rather than a constant in a runner script: the two geometries
address different cache entries and appear as different rows.

\subsection{Input resolution}
\Cref{tab:geometry} varies input size while holding the encoder, the head and the protocol
fixed. Three regularities hold across both C-RADIOv4 capacities.

\begin{table*}[t]
\centering
\caption{Input geometry, ArcFace mAP. Rows 1--3 and 4--6 vary resolution within one
checkpoint; rows 7--8 vary only the resize convention. \emph{native} feeds each image at its
own size, snapped to the patch grid and capped at 512 pixels.}
\label{tab:geometry}
\scriptsize
\setlength{\tabcolsep}{3pt}
% generated by tools/gen_tables.py from results/runs — edit the runs, not this
\begin{tabular}{lllrrrrrrr}
\toprule
encoder & input & fit & Market & MSMT17 & CUHK03 & Occ-REID & CCVID & VRIC & VRAI \\
\midrule
C-RADIOv4-H & 224$\times$224 & squash & 0.709 & 0.165 & 0.153 & 0.646 & 0.609 & 0.191 & 0.228 \\
C-RADIOv4-H & 256$\times$128 & squash & \textbf{0.718} & \textbf{0.180} & \textbf{0.163} & 0.622 & 0.612 & 0.185 & 0.187 \\
C-RADIOv4-H & \textit{native} & squash & 0.559 & 0.127 & 0.119 & 0.460 & 0.589 & 0.106 & \textbf{0.244} \\
\midrule
C-RADIOv4-SO400M & 224$\times$224 & squash & 0.696 & 0.162 & 0.140 & \textbf{0.651} & 0.625 & \textbf{0.197} & 0.207 \\
C-RADIOv4-SO400M & 256$\times$128 & squash & 0.718 & 0.175 & 0.149 & 0.647 & 0.581 & 0.177 & 0.179 \\
C-RADIOv4-SO400M & \textit{native} & squash & 0.573 & 0.136 & 0.119 & 0.481 & 0.612 & 0.097 & 0.225 \\
\midrule
CLIP-B/16 (squash) & 224$\times$224 & squash & 0.289 & 0.027 & 0.017 & 0.414 & \textbf{0.795} & 0.035 & 0.039 \\
CLIP-B/16 (crop) & 224$\times$224 & crop & 0.174 & 0.016 & 0.018 & 0.437 & 0.630 & 0.049 & 0.035 \\
\bottomrule
\end{tabular}

\end{table*}

First, \textbf{$256\times128$ wins on Market-1501 and CUHK03-NP and loses on both vehicle
sets}, at both C-RADIOv4 capacities, while also running about 1.6$\times$ faster; where it
wins it is therefore the better purchase on both axes at once.

Second, the tempting explanation for that --- that a $2{:}1$ input matches a person crop's
aspect ratio --- is \textbf{false}, and this study contains its own counterexample.
Market-1501 and Occluded-REID ship images that are every one of them exactly $64\times128$:
identical median size, identical $2.00$ aspect, coefficient of variation $0.00$ in image area.
With the same head they prefer \emph{opposite} input sizes --- $256\times128$ on Market
($0.718$ against $0.709$) and $224\times224$ on Occluded-REID ($0.646$ against $0.622$, about
five times that set's seed noise). CUHK03-NP, at a much taller $2.97$ aspect, prefers $2{:}1$
by \emph{more} than Market does rather than less. A second candidate explanation, token count,
fails on the same rows. We report the regularity and explicitly do not claim a mechanism for
it; ``$2{:}1$ usually wins on person crops'' is a useful default with two measured exceptions
inside this table, not a law.

Third, \textbf{\emph{native} resolution wins on exactly one dataset, VRAI, and loses
everywhere else --- including on the other vehicle set}. That last detail is the one that
identifies the variable. If the split were person-versus-vehicle, native should also win on
VRIC; instead it collapses there, from $0.191$ to $0.106$. What separates the two is crop
size. VRAI's median crop is $295\times202$, inside the range C-RADIOv4 was trained across;
VRIC's crops, and every Market-1501 and Occluded-REID crop, are far below it. Indeed on the two
person sets where every file is exactly $64\times128$, ``native'' is not native at all: it is a
fixed 32-token resolution wearing a misleading label, well under the roughly 128-pixel floor
the encoder saw in training.

The variable that moved is therefore resolution \emph{relative to the encoder's training
distribution}, not resampling and not object class. Upsampling a small crop to $224$ is not a
distortion to be avoided; it is how the crop reaches a size the encoder can read. Native is
also far cheaper on small crops --- roughly 4$\times$ the throughput --- so it remains a real
option wherever a quarter of the mAP is an acceptable price.

Two further cells resist the story and we report them rather than smooth them: CLIP's
Occluded-REID and VRIC rows prefer the centre-crop geometry, which every other column in the
table penalises. Both are among the weakest scores in the study, and we have not established
whether that is the reason.

\subsection{Why this belongs in the paper rather than in an appendix}
Take the two variables together. Resize geometry is worth up to a large fraction of a probed
score; resolution is worth on the order of a fifth to a quarter of one. Several of the backbone
comparisons in \cref{tab:ablation} are decided by less than either. A frozen-encoder comparison
that does not pin both is therefore not merely imprecise --- it is capable of reporting the
reverse of the truth, and it will do so with a table that looks exactly like a correct one.
The remedy is cheap and we recommend it as a reporting standard: state the resize convention
and the input size beside every frozen-encoder number, and treat any comparison that omits
them as under-specified.

\begin{table}[t]
\centering
\caption{Extraction throughput on this study's single RTX 2070 Max-Q, Market-1501, frozen
features. A throughput figure measures a machine and a moment, not a model --- one cell here
first recorded 2.2 img/s against 12.8 for the same encoder elsewhere because
another process held part of the card; re-extracting on a quiet card gave 12.3, and the
metrics never moved.}
\label{tab:cost}
\scriptsize
% generated by tools/gen_tables.py from results/runs — edit the runs, not this
\begin{tabular}{lllr}
\toprule
encoder & input & params (M) & img/s \\
\midrule
C-RADIOv4-H & 224$\times$224 & 653 & 12.3 \\
C-RADIOv4-H & 256$\times$128 & 653 & 20.8 \\
C-RADIOv4-H & \textit{native} & 653 & 53.9 \\
C-RADIOv4-SO400M & 224$\times$224 & 431 & 21.3 \\
C-RADIOv4-SO400M & 256$\times$128 & 431 & 31.4 \\
C-RADIOv4-SO400M & \textit{native} & 431 & 85.9 \\
SigLIP2-g & 256$\times$256 & 1163 & 5.8 \\
SigLIP2-SO400M & 224$\times$224 & 428 & 15.6 \\
DINOv3-H+ & 224$\times$224 & 840 & 9.5 \\
DINOv3-L & 224$\times$224 & 303 & 25.5 \\
CLIP-B/16 (squash) & 224$\times$224 & 87 & 81.8 \\
CLIP-B/16 (crop) & 224$\times$224 & 87 & 33.9 \\
\bottomrule
\end{tabular}

\end{table}
